\scanpage{72}

% top across sheet, diagonally cancelled comparison; heavily crossed first pair
\diagram{\deleted{\resizebox{\linewidth}{!}{\begin{tikzcd}[ampersand replacement=\&,column sep=small,row sep=small]
\deleted{\unsure{C^*(X,\Omega^*)}}\arrow[r] \& \deleted{\unsure{\mathbb R\Gamma(X,\Omega^*)}} \& \mathbb R^*\Gamma(X,k)\arrow[l,"\substack{\varepsilon_{\Omega^*}\\\sim}"']\arrow[r,"\substack{\varepsilon_{S_X^*}\\\sim}"]\arrow[d,"\sim"]\arrow[dr,dashed,"\varepsilon_{S_X^{\prime *}}"] \& \mathbb R^*\Gamma(X,S_X^*)\arrow[d,"\sim"] \& \Gamma(X,S^*_{X,\deleted{\unsure{k}}})\arrow[l,"\substack{(S_X^*\ \text{\unsure{étant} par Cartan}\\\text{hom. fin)}\\\sim}"']\arrow[d,"\sim"] \& {}=C^*(S(X)_{\unsure{*|0}},k)\arrow[d,"\sim"]\\
\& \& \mathbb R^*\Gamma(X,\Omega^*)\arrow[r,"\sim"] \& \mathbb R^*\Gamma(X,S_X^{\prime *}) \& \Gamma(X,S_X^{\prime *})\arrow[l,"\sim"'] \& {}=C^*(S(X)'_*,k)\\
\& \& C^*(X,\Omega^*)\arrow[u] \& \& \& {}
\end{tikzcd}}}}

% upper middle, NB introducing the augmentation square
\uline{NB} \unsure{l'hyp.|les hom|la donnée} (4) joint avec $\unsure{K_*|k_*}$ définit un hom

% upper middle right, four-object square embedded in the note
\diagram{\begin{tikzcd}[ampersand replacement=\&,row sep=small]
\Omega^*\arrow[r,"\inserted{\text{par }K}"] \& S_X^{\prime *}\\
k\arrow[u]\arrow[r,dash] \& S_X^*\arrow[u]
\end{tikzcd}}

% upper middle, sentence after square
compatible avec les augmentations, d'où carré comm. \deleted{\unsure{(i) on|(i) ou}} (d'isom dans $D(X,S)$) donnant lieu à un cube commutatif

% center, cube with dashed back edges, Cartan and hypothesis labels
\diagram{\resizebox{\linewidth}{!}{\begin{tikzcd}[ampersand replacement=\&,column sep=small,row sep=large]
C^*(X,\Omega^*)\arrow[rr]\arrow[dd,"\deleted{\unsure{\sim|\wr}}"'] \& \& C^*(X,S_X^{\prime *})\overset{\mathrm{dfn}}{=}C(S'_X,k)\arrow[dd,"\sim"] \& {}\\
\& {[C(X,k)]}\arrow[ul,dashed]\arrow[ul,shift left=1.3ex,"\deleted{\unsure{\leftrightarrow}}" description]\arrow[rr,dashed]\arrow[dd,dashed] \& \& C^*(X,S_X^*)=C(\unsure{S_X|S_x},k)\arrow[ul,"\substack{\text{hyp. sur }S'_X\\\sim}"']\arrow[dd, "{\substack{\sim\ (\text{Cartan, }S_X^*\text{ est}\\\text{``homot.\textsuperscript{t} fin'')}}}"]\\
\mathbb R^*\Gamma(X,\Omega^*)\arrow[rr,"\sim"] \& \& \mathbb R^*\Gamma(X,S_X^{\prime *}) \& {}\\
\& \mathbb R^*\Gamma(X,k)\arrow[ul,"\sim"] \& \& \mathbb R^*\Gamma(X,S_X^*)\arrow[ul,"\sim"]\arrow[ll,"\sim"']
\end{tikzcd}}}

% right margin at cube, written vertically
% margin note written sideways
\margin{\deleted{$\longleftarrow\mathbb R^*\Gamma(X,S^*_{X,k})$}}

% below cube, last four lines
La base du cube est formée d'isom, car transformée par $\mathbb R^*\Gamma$ d'un carré d'iso dans une cat. dérivée ; la face \unsure{devant|dev$^{\mathrm t}$|du haut} est formée d'iso, car on le sait déjà pour \unsure{deux entre eux|deux autres}, donc aussi pour le quatrième qui reste.

% lower left
% archive sheet number: 71
